% =====================================================================
%  06-conclusiones.tex
% =====================================================================
\chapter{Conclusiones}\label{cap:conclusiones}

Este capítulo presenta las conclusiones del corte técnico del 12 de septiembre
de 2026, sustentadas en la implementación y las pruebas de los
\cref{cap:implementacion,cap:pruebas}. El avance permite evaluar el núcleo
criptográfico y el flujo documental autenticado, pero no declarar terminado
el conjunto de objetivos funcionales del Trabajo Terminal. La gestión procesal
y la evaluación con usuarios conservan trabajo pendiente.

\section{Resultados obtenidos}\label{sec:resultados-obtenidos}

El prototipo integra cálculo de resúmenes, cifrado autenticado, firma digital,
validación de certificados y revocación, sellado temporal y exportación de
paquetes de evidencia. La comprobación con OpenSSL permite examinar documento,
firma, certificado y sello sin depender de la aplicación que los produjo.
La autoridad temporal utilizada es local: demuestra la interoperabilidad del
protocolo RFC 3161 bajo las anclas del ensayo, sin sustituir una autoridad
independiente ni acreditar por sí misma los efectos jurídicos de la evidencia.

El flujo HTTP incorpora expedientes y asignaciones persistidas. Cada operación
documental comprueba la sesión y revalida el rol activo, la pertenencia vigente
y la asociación entre documento y expediente al confirmar su acceso.
Owner conserva acceso global a expedientes; Litigante y Paralegal requieren
asignación, Paralegal no puede sellar y Cliente mantiene denegadas las
operaciones documentales. El conocimiento de un identificador no basta para
acceder al recurso.

La persistencia PostgreSQL resuelve la separación entre cambio documental y
registro de éxito: ambos se confirman o revierten dentro de la misma
transacción. Expedientes, asignaciones y mutaciones de identidad prolongan
la misma cadena de auditoría. El trabajo criptográfico y la comunicación con
la TSA se preparan fuera de esa transacción; los commits se ordenan mediante
un bloqueo compartido por los escritores. La asociación documental y la
evidencia ya sellada permanecen inmutables.

La evaluación reproducida aprobó 519 pruebas, sin fallos y con una prueba
externa ignorada; la cobertura fue de 8026 de 8845 líneas, equivalente al
90.7\,\%. El ensayo de dos servidores produjo un sellado aceptado y un
conflicto, con un único evento y paquetes ZIP idénticos. La migración y
restauración utilizaron cuatro documentos, tres sellados, y conservaron las
46 entradas históricas de auditoría y sus hashes. La comparación del estado
restaurado y de los paquetes, junto con la verificación OpenSSL, aporta evidencia
de recuperación para ese conjunto controlado. Los datos fueron desechables;
no se trató de una migración de expedientes reales del despacho.

\section{Contribuciones}\label{sec:contribuciones}

La primera contribución es una separación verificable entre reglas de dominio,
casos de uso y adaptadores. Los puertos permiten compartir la preparación
criptográfica entre la CLI y el flujo HTTP sin introducir persistencia ni
llamadas externas dentro del dominio. La extracción de esta preparación
facilita probar autorización y fallos de almacenamiento sin regenerar evidencia.

La segunda contribución es una frontera transaccional que conecta autorización,
mutación e historia. Las pruebas de revocación y concurrencia hacen explícito
cuándo una operación puede confirmarse, mientras que los permisos del rol
operativo restringen la reescritura de documentos y eventos. Este resultado
acota una propiedad concreta del backend; no presupone que un administrador
de la base de datos quede fuera de la base de confianza.

La tercera contribución es un procedimiento de migración que conserva los
bytes criptográficos y reconcilia las fuentes con el destino. La inspección sin
escrituras, el mapa explícito de asociaciones, las barreras de escritura y los
recibos permiten distinguir fallos anteriores y posteriores al commit. La guía
operativa y el ensayo de restauración complementan las pruebas funcionales con
un procedimiento reproducible de recuperación.

\section{Limitaciones}\label{sec:limitaciones}

La validación corresponde a un entorno controlado. No se midieron carga de
producción, disponibilidad prolongada ni usabilidad con personal del despacho.
Siguen pendientes TLS, administración de conexiones, plazos operativos,
apagado ordenado y revisión del modelo de amenaza de la firma RSA expuesta
por HTTP. El bloqueo global simplifica el orden de la auditoría a costa de
serializar las confirmaciones; su costo debe medirse antes de aumentar la carga.

Redis no participa en una transacción distribuida con PostgreSQL. Ante un fallo
de auditoría se intenta retirar el desafío o la sesión recién creados y no se
entrega su token; si también falla la limpieza, la clave puede permanecer hasta
su expiración. Asimismo, una respuesta cuyo acceso ya se confirmó puede terminar
de transmitirse después de una revocación posterior. Las pruebas fijan ese orden;
no afirman la cancelación retroactiva de respuestas.

Una cadena de hashes válida permite detectar alteraciones interiores, pero no
prueba por sí sola que su extremo no haya sido truncado. El anclaje externo de
la cabeza sigue pendiente. La migración comprueba cadena y CRL del firmante en
el instante del sello; la validación de confianza de la TSA mantiene la política
actual de OpenSSL y puede rechazar una autoridad expirada al verificarla.
Preservar la evidencia permite examinarla posteriormente, pero no garantiza su
aceptación indefinida ni una decisión jurídica determinada.

La carga crea únicamente la versión uno. No existen todavía consultas de
listado, búsqueda e historial documental, ni cambio de contenido o reasignación
del documento a otro expediente. Las asignaciones de usuarios representan
acceso al sistema, no participantes procesales. Tampoco se han implementado
las audiencias, plazos, alertas o la interfaz de usuario completa.

\section{Trabajo futuro}\label{sec:trabajo-futuro}

La siguiente ampliación debe ofrecer listado, detalle, búsqueda e historial
documental con autorización por expediente. Cada versión deberá preservar su
identidad, contenido cifrado y evidencia histórica, con pruebas tanto de acceso
permitido como de rechazo. La política documental de Cliente requiere un diseño
y validación explícitos antes de ampliar sus permisos.

Sobre esa base corresponde modelar participantes procesales, audiencias y plazos,
sin confundirlos con las asignaciones de acceso ya existentes. Después se debe
integrar la interfaz de login y MFA, navegación, carga, sellado, verificación y
descarga contra el contrato HTTP, y evaluar los flujos con usuarios del despacho.
El cronograma de diseño conserva su carácter de planificación; estos resultados
no permiten dar por cumplidas sus actividades de usabilidad y gestión procesal.

Antes de un despliegue público se deben ensayar los límites operativos señalados,
restauraciones con volúmenes representativos y fallos de servicios externos. La
contratación y validación de un prestador autorizado y el anclaje independiente
de auditoría son extensiones distintas de la demostración local. Cada avance
funcional deberá actualizar el contrato, las decisiones, las mediciones y los
apartados correspondientes de este reporte, manteniendo separados los resultados
reproducidos, las hipótesis y el trabajo pendiente.
